\color{unidarkgreen}{\subsection{Токовая защита нулевой последовательности стороны ВН (ТЗНП ВН)}}
\color{black}

\subsubsection{Общие сведения}

\begin{enumerate}[label=\arabic{section}.\arabic{subsection}.\arabic{subsubsection}.\arabic*,  align=left, labelsep=4pt, leftmargin=0pt, itemindent=57pt]

\item
Токовая защита от замыканий на землю на стороне ВН используется при наличии питания с других
сторон трансформатора. Защита предназначена для резервирования отключения замыкания на землю
на шинах (и линиях) стороны ВН трансформатора, а также для резервирования основных защит
трансформатора.

\item
Защита подключается к ТТ, установленным со стороны обмоток ВН силового трансформатора.

\item
В состав ТЗНП ВН входит одна ненаправленная ступень с контролем тока нулевой
последовательности высокой стороны трансформатора. Порог срабатывания ступени защиты задается
уставкой \textbf{«3I0ср ТЗНП ВН»}, а выдержка времени срабатывания ступени задается уставкой \textbf{«Тср ТЗНП
ВН»}. Ввод в работу ТЗНП ВН осуществляется с помощью уставки \textbf{«ТЗНП ВН»}. Оперативный ввод/вывод
ТЗНП ВН производится с помощью БУ (блока управления функциями) «Вывод ТЗНП ВН». Блокировка по
2 гармонике выводится в логике ТЗНП накладкой \textbf{«Режим БНТ ТЗНП ВН»}.

\item
Для обеспечения несрабатывания ТЗНП ВН при броске тока намагничивания используется
блокировка по содержанию 2 гармоники в токе нулевой последовательности стороны ВН. Блокировка
основана на контроле отношения действующего значения второй гармонике к первой тока нулевой
последовательности. Данная блокировка вводиться в работу с помощью уставки \textbf{«БНТ ТЗНП ВН»}.
Имеется возможность регулировать порог срабатывания блокировки, с помощью уставки \textbf{«3I02г/3I01г
ТЗНП ВН»}. Блокировка ТЗНП ВН по 2 гармонике выводится из действия при величине тока нулевой
последовательности ниже 0,04Iном. Также блокировка ТЗНП ВН по 2 гармонике выводится из действия
при достижении величины действующего значения тока нулевой последовательности выше уставки \textbf{«Iср
БНТ ТЗНП ВН»}.

\begin{figure}[H]
\centering
\includegraphics[width=0.8\textwidth,height=0.8\textheight,keepaspectratio]{\fbpath/021.0141 ТЗНП ВН/img/TZNP_VN.pdf}
\captionof{figure}{Функциональный блок ТЗНП ВН}\label{pic:TZNP_VN}
\end{figure}


\small
\begin{longtable}{|>{\centering\arraybackslash}m{5.3cm}|>{\centering\arraybackslash}m{4.2cm}|>{\centering\arraybackslash}m{1cm}|>{\centering\arraybackslash}m{3.3cm}|>{\centering\arraybackslash}m{1.8cm}|}
% Первая страница: полный заголовок
\caption{Перечень уставочных параметров ТЗНП ВН\hfill\vspace{-0.5\baselineskip}}\label{desc:tbl1}\\
\hline
\rowcolor{gray!30}
Наименование уставки & Диапазон & Шаг & Значение по умолчанию & Единица измерения \\
\hline
\endfirsthead

% Промежуточные страницы: «Продолжение таблицы»
\caption*{\hspace{3pt}\emph{Продолжение таблицы \ref{desc:tbl1}}\hfill\ }\vspace{-0.5\baselineskip} \\
\hline
\rowcolor{gray!30}
Наименование уставки & Диапазон & Шаг & Значение по умолчанию & Единица измерения \\
\hline
\endhead

%===m>ДТЗ|-
%===m
\end{longtable}
\normalsize	

\end{enumerate}